\section{Masked-data inference: brief primer}
\label{sec:background}

We summarize the masked-data identifiability framework
\citep{towell2026masked,heitjan1991ignorability}. Readers familiar
with this work or with \citet{towell2026scrnacoarsening} can skim to
\cref{sec:translation}.

Consider $m$ independent components in series. Each component $j$
has a parametric distribution $f_j(\cdot;\theta_j)$ for its lifetime
$T_{ij}$ in system $i$. The system fails at $T_i = \min_j T_{ij}$
with cause $K_i = \arg\min_j T_{ij}$. We observe $T_i$ but not $K_i$
directly; instead we observe a \emph{candidate set} $C_i \subseteq
\{1,\ldots,m\}$ thought to contain the failed component.

The joint likelihood factorizes as
\begin{equation}
  f(t, k, c; \theta)
  = h_k(t; \theta_k) \prod_{\ell} R_\ell(t; \theta_\ell)
    \cdot \Prob_\theta\{C_i = c \mid T_i = t, K_i = k\},
  \label{eq:bg-joint}
\end{equation}
where $h_\ell$ is the component-$\ell$ hazard and $R_\ell$ its
survival function. The third factor is the unknown masking mechanism.

Three coarsening conditions allow the masking factor to drop from
the score:

\begin{condition}[Support, C1]
\label{cond:c1}
$\Prob\{K_i \in C_i\} = 1$.
\end{condition}

\begin{condition}[Symmetry, C2]
\label{cond:c2}
$\Prob_\theta\{C_i=c \mid T_i=t, K_i=j\} = \Prob_\theta\{C_i=c \mid
T_i=t, K_i=k\}$ for all $j,k \in c$.
\end{condition}

\begin{condition}[Parameter independence, C3]
\label{cond:c3}
$\Prob_\theta\{C_i=c \mid T_i, K_i\}$ does not depend on $\theta$.
\end{condition}

Under all three, the likelihood collapses to
\begin{equation}
  L_i(\theta) \propto \left[\prod_\ell R_\ell(t_i;\theta_\ell)\right]
  \cdot \left[\sum_{j \in c_i} h_j(t_i; \theta_j)\right].
  \label{eq:bg-c1c2c3}
\end{equation}

\citet[Theorem 8]{towell2026masked} establish identifiability:

\begin{theorem}[Identifiability under C1--C2--C3]
\label{thm:bg-id}
Under \cref{cond:c1,cond:c2,cond:c3} and standard regularity, full
column rank $m$ of the augmented candidate-set matrix $\tilde{C}
\in \{0,1\}^{n \times m}$ is \emph{sufficient} for the parameter
vector $\theta$ to be identifiable. Separability of the masking
mechanism is necessary but not in general sufficient: a separating
mechanism can leave $\tilde{C}$ rank-deficient (for example $m=4$
with candidate sets $\{1,2\},\{3,4\},\{1,3\},\{2,4\}$), in which
case $\theta$ is not identifiable without further functional
restrictions on the component family.
\end{theorem}

\emph{Singleton} candidate sets ($|c_i|=1$) are critical: each
isolates a single component contribution in
\eqref{eq:bg-c1c2c3} and adds an identifying row to $\tilde{C}$.

\citet{towell2026scrnacoarsening} ported this framework to
single-cell RNA sequencing (scRNA-seq) zero inflation. We now port
it to spatial transcriptomics, where the candidate-set structure is
richer: instead of a single bit (zero vs. nonzero) per gene, each
spot carries a probability distribution over cell types.
